\begin{lemma}[Step 1 exit-leg global minimum]
Let $T$ be a positive-rung ETT of rung number $n$, constructed from an
arbitrary certified input coloring and with generating coloring
$\varphi_n$.  Assume that $T$ has maximum property, that $\varphi_n$ is
proper and leaves exactly the distinguished edge $e_0$ uncolored, and
fix a hierarchy base.  Let
\[
 T_n=T_{n,0}\subset T_{n,1}\subset\cdots\subset
 T_{n,q}\subset T_{n,q+1}=T
\]
be a good hierarchy $H$ over that base, with its actual family of
two-color sets $\Gamma^j_\eta$.

Assume, uniformly over every certified input coloring and generating
coloring, that every positive-rung ETT $U$ of rung number $s<n$ with
maximum property satisfies all six lower-rung assertions:
\begin{enumerate}
\item $V(U)$ is elementary under its generating coloring;
\item under every proper coloring leaving exactly $e_0$ uncolored and
stable on the final ladder closure, that closure is interchangeable for
two colors whenever at least one is missing there;
\item stability descends at every PE support to the prefix before that
support and to its first-appearance ladder stage, and the stored pairs at
distinct PE slots with distinct supports and the same first-appearance
index are disjoint;
\item at a terminal PE, the stored two-color component through the
support meets the terminal ladder tree only at the support;
\item every nonminimum endpoint of a defective boundary color precedes
each later PE support or SE extension vertex allowed by its
first-appearance index; and
\item every proper terminal stable coloring leaving exactly $e_0$
uncolored is a mod-coloring of $U$.
\end{enumerate}

Assume also the induction hypothesis (6.1) in the following explicit
form.  For every arbitrary certified input coloring and generating
coloring $\psi$, every ETT $U$ of the same rung number $n$ as $T$ which
has maximum property, and every good hierarchy $K$ of $U$ with strictly
fewer than $q$ dividers, $V(U)$ is elementary under $\psi$.  Finally
assume the large-branch estimate
\[
              |\overline{\varphi_n}(T_n)|\geq 2n+11.       \tag{L}
\]

Let $1\leq p\leq q$, let
$S=T_{n,p}$ (formally, for $i=p-1$ this is
`$H.\mathit{stage}(i.\mathit{succ}.\mathit{castSucc})$'), and let
$\alpha,\beta\in[k]$ be distinct.  Assume
$\alpha\in\overline{\varphi_n}(S)$, the singleton $\{\alpha\}$ is closed
on $S$ under $\varphi_n$, and $\alpha,\beta$ are not
$S$-interchangeable.

Then there exist an arbitrary certified input, a proper
generating coloring $\psi$ leaving exactly $e_0$ uncolored, and a
same-rung ETT $U$ on an exact terminal tree $T'$ with maximum property,
together with a base and a good hierarchy $K$, such that
$d=K.\mathit{dividerCount}<q$.  This presentation has its own large
estimate, its separate terminal minus-closed certificate, and
$V(T')$ is elementary under $\psi$.

There exist distinct colors $a,b$, a color $\theta$, two exact
support-disjoint $(T',\psi,\{a,b\})$-exit paths $P_1,P_2$ with distinct
outside and tree ends $a_1,b_1,a_2,b_2$ ordered by $b_1\prec b_2$, and
individual $b$-colored boundary-edge identities $f_1,f_2$ on the paths
incident with $b_1,b_2$.  The color $a$ is missing and singleton-closed
on $T'$.  Writing $A=K.\mathit{starred}(d)$, the color $\theta$ is
different from $a,b$, is missing and singleton-closed on $A$, and is
singleton-closed on $T'$.

This presentation is primary-minimal over the full replacement family:
under every proper exactly-one-uncolored stable coloring, no
corresponding same-rung maximum ETT with its own large estimate,
terminal certificate, and good hierarchy with fewer than $d$ dividers
admits two such exit legs.  Moreover, either $b$ is missing on $T'$, or
$b$ is not missing and there is a third exact exit path $P_3$, pairwise
support-disjoint from the first two, with a third outside and tree end,
$b_2\prec b_3$, and a retained $b$-colored boundary identity $f_3$ on
$P_3$.  In the latter alternative the chosen triple is secondary-minimal
over all corresponding same-divider presentations: none has second
color missing, and none whose second color is present has a qualifying
triple with total exact walk length smaller than
$|P_1|+|P_2|+|P_3|$.  Every comparison ranges over arbitrary certified
replacement inputs and retains individual parallel-edge identities.
\end{lemma}
\begin{proof}
Write $\varphi=\varphi_n$.  We first isolate the presentations on which
the induction will range.  For each $1\leq r\leq q$, apply
\noderef{GoodHierarchyPrefixRepackage} to the given maximum ETT, base,
hierarchy $H$, and good family $\Gamma$.  It returns a maximum ETT on the
exact ordered tree $T_{n,r}$ with the same certified series and individual
connecting-edge identities, the same base data, and a good hierarchy whose
divider count is $r-1$.  Its retained stages and divider positions are
exactly the old ones, and its stored sets and paper Gamma values agree at
every retained index.  The theorem also returns, separately, the old
`minusClosed' field at index $r-1$, namely that the new terminal tree is
$(\bigcup_{\eta\in D_{n,r}}\Gamma^{r-1}_\eta)^-$-closed; it does not
incorrectly claim this terminal certificate as one of the shorter good
hierarchy's internal fields.
The same-rung induction hypothesis (6.1), applied to this exact
truncated ETT and hierarchy, gives
\[
             V(T_{n,r})\text{ is elementary under }\varphi.
                                                               \tag{1}
\]
This construction works for arbitrary certified input and generating
colorings, which is essential when we repeat it after a recoloring.

We shall also truncate at a vertex $b$ lying strictly after
$A=T_{n,r-1}^{*}$ and no later than $T_{n,r}$.  Use
\noderef{TreeSequencePrefix} to take the ordered prefix through $b$, or
the prefix $B$ immediately before $b$.  Every retained vertex and edge
is the same indexed identity by \noderef{TreeSequenceInitialSegment}.
There are three genuinely different cases.

If $A$ is a strict initial segment of $B$, keep the stages through
$T_{n,r-1}$ and use $B$ as terminal.  The last strictness field of
\noderef{TashkinovHierarchy} is now exactly $A\subsetneq B$, so this is
the valid $r-1$-divider presentation.  If instead $B=A$ and $r>1$, do
not use that presentation: retain only the stages through
$T_{n,r-2}$ and take $T_{n,r-1}=A$ as terminal.  This is the valid
$r-2$-divider presentation, with its terminal minus-closed certificate
given by the old field at $r-2$.  In either case a recoloring on a full
component disjoint from $B$ is an \noderef{OutsideKempeRecoloring} of
this displayed ETT.  Once its two colors have been proved distinct,
\noderef{HierarchyLEMStable} supplies properness, exact
one-uncoloredness, the corresponding ETT, maximum property, the same
ordered stages, stored sets and $\Gamma$-pairs, and the transported
terminal minus-closed certificate.

The remaining case is $r=1$ and $B=A=T_{n,0}^{*}$.  We do not present
the starred base as an ETT.  The component is disjoint from $A$ and hence
from $T_n$.  Thus \noderef{OutsideKempeRecoloringStable} supplies the
proper, one-uncolored, terminal-ladder stability data.  Lower assertion
(vi), instantiated by \noderef{Td}, makes the changed coloring a
\noderef{ModColoring}; the zero-rung stability and interchangeability
endpoint is \noderef{TechnicalTheoremBase}.  This is the separate base
route used below.

For the one-edge enlargement from $B$ to the prefix $C$ through $b$,
the old final certified closure in
\noderef{TashkinovSeriesStep} and \noderef{TashkinovClosure} supplies the
exact finite augmentation sequence.  Because the recoloring fixes every
edge identity and every missing-color witness on $B$,
each retained augmentation still satisfies
\noderef{TashkinovAugmentation}; the universal clause of
\noderef{ModColoring}, obtained from lower assertion (vi) via
\noderef{Td}, then supplies the corresponding ETT on the enlarged
prefix.  \noderef{MaximumPropertyTransport} transfers maximum property.
When $B=A$ and $r>1$, adjoining $b$ restores the $r-1$-divider
presentation: the last strict containment is the exact one-edge
containment $A\subsetneq C$.  Pair size, disjointness, unique refresh,
parallel base, and every old minus-closed field are copied; final
avoidance is checked on that exact new edge.  When $r=1$, the same check
constructs the zero-divider hierarchy with base stage $T_n$, terminal
$C$, and core inside $C$.  Thus no hierarchy is ever asserted with its
last retained stage equal to its terminal tree, and no endpoint pair is
substituted for a graph-edge identity.

For a displayed tree $U$, a proper coloring $\psi$, and distinct colors
$a,b$, define the following local proof predicate:
\[
 \mathsf{LegBad}(U,\psi,a,b)
\]
means that $a\in\overline\psi(U)$, the singleton $\{a\}$ is closed on
$U$ in the sense of \noderef{TreeColorSetClosed} (so
\noderef{BoundaryColorEdges} contains no $a$-colored boundary identity),
and there are two \noderef{TreeExitPath} witnesses for the colors
$a,b$ whose outside endpoints are distinct and whose node sets as walks
in \noderef{TwoColorIncidenceGraph} are disjoint.  In particular, the
two witnesses have distinct tree ends and distinct first boundary-edge
identities.  This predicate says nothing about whether their full
two-color components are distinct.

We record two consequences which will be used repeatedly.  First, let
$U$ be the terminal tree of any same-rung maximum ETT with a good
hierarchy having $d<q$ dividers.  The same-rung hypothesis (6.1) makes
$V(U)$ elementary.  If $a$ is missing and singleton-closed on $U$ and
$a,b$ are not $U$-interchangeable, the negation of
\noderef{TreeInterchangeable}, with
\noderef{TwoColorPathComponent} unfolded, gives two distinct paper path
components meeting $U$; one may initially be a singleton.  Their full
node sets are disjoint by \noderef{TwoColorComponentsDisjoint}.

If $b$ is not missing on $U$, a singleton component meeting $U$ is
impossible: its unique tree vertex would have no incident edge of either
path color and hence would miss $b$.  Thus both components contain edge
nodes and satisfy \noderef{TwoColorNontrivialPathComponent}.  Each has a
simple spanning path by \noderef{TwoColorComponentSpanningPath};
properness and \noderef{TwoColorIncidenceDegreeAtMostTwo} show that its
two ends are vertex-nodes missing $a$ or $b$.  Only the unique
$a$-missing vertex can be an end in $U$, so at least three of the four
ends are outside $U$.

Now suppose that $b$ is missing on $U$.  If both components are
nontrivial, the same spanning-path argument applies, and elementariness
permits at most one of their four ends in $U$ to miss $a$ and at most one
to miss $b$; hence at least two ends are outside.  If, after relabeling,
the first component is a singleton at $x\in U$, then both $a$ and $b$
are missing at $x$.  Elementariness makes the other component nontrivial
and prevents either end of its spanning path from lying in $U$, since
every such end misses $a$ or $b$.  Its two ends therefore supply the two
required outside ends.  There cannot be two distinct singleton
components, since their two tree vertices would repeat both missing
colors.

From each selected outside end, follow its component's spanning path only
as far as its first vertex-node in $U$.  The component meets $U$, so this
portion exists.  Its outside end misses one of the two colors, it has at
least one edge, and its only tree vertex is its last vertex; hence it is a
\noderef{TreeExitPath}.  Portions in different full components are
disjoint.  Two portions selected from opposite ends of one spanning path
are also disjoint: if their first contacts were equal, their distinct
last edge-nodes would both be $b$-colored, because the closed color $a$
has no boundary identity, and properness would fail at that contact.
Selecting the required ends therefore yields pairwise disjoint
first-contact portions.  In particular, failure of interchangeability
implies $\mathsf{LegBad}(U,\psi,a,b)$, and when $b$ is not missing it
supplies a qualifying triple of disjoint exit legs with distinct outside
endpoints.

Second, the same qualifying triple exists whenever
$\mathsf{LegBad}(U,\psi,a,b)$ holds and $b$ is not missing, even though
the two given legs may belong to one full component.  The outside end of
each given leg is a degree-one vertex-node: it misses one path color and
is incident with the other along the nontrivial leg.  If the two ends
belong to distinct full components, those components have four ends and
at most the unique $a$-missing vertex can lie in $U$, so three ends are
outside.  If they belong to one full component, they are its two ends.
Let $x$ be the unique vertex of $U$ missing $a$.  Since $b$ is present at
$x$, the full component through $x$ is nontrivial and has degree-one end
$x$; it differs from the component already having the two outside ends.
Its other end cannot lie in $U$, since no vertex of $U$ misses $b$ and
no second vertex there misses $a$.  It is therefore a third outside end.
Taking first-contact portions as in the preceding paragraph gives the
claimed disjoint triple component by component.  No parity assertion
about raw cut edges is used.

For the zero-divider endpoint, apply \noderef{Step1ZeroDividerAnchor} to the exact hierarchy core and the two retained exit legs.  Its non-PE and PE branches prove that no admissible zero-divider presentation exists.  This supplies the anchor required below.

We now define the well-founded comparison domain.  An
\emph{admissible leg presentation} consists of an arbitrary certified
input coloring, a proper generating coloring $\psi$ leaving exactly
$e_0$ uncolored, a corresponding same-rung ETT $U$ with maximum
property, a hierarchy base and a good hierarchy $K$ with $d<q$
dividers, the fixed lower six-assertion package and (6.1), the analogue
of (L) for its own terminal ladder tree, colors $a\ne b$ satisfying
$\mathsf{LegBad}(U,\psi,a,b)$, and the separate terminal certificate
\[
 U\text{ is }
 \left(\bigcup_{\eta\in D^K_{n,d+1}}\Gamma^d_\eta\right)^-
 \text{-closed under }\psi.                                  \tag{C}
\]
Here the displayed $D^K_{n,d+1}$ is the terminal expression obtained by
unfolding \noderef{TashkinovStoredColors} and
\noderef{HierarchyStoredColorsAt}; it is not an extra good-hierarchy
field.

Give a presentation primary coordinate $d$.  If $b$ is missing on $U$,
give it secondary coordinate $0$.  Otherwise the second consequence of
$\mathsf{LegBad}$ above supplies qualifying triples of pairwise disjoint
exit legs with distinct outside ends; give it secondary coordinate one
plus the least total length of such a triple.  The sets minimized are
nonempty finite-length subsets of $\mathbb N$, and lexicographic order on
$\mathbb N\times\mathbb N$ is well founded.

The assumed failure of interchangeability makes the comparison class
nonempty.  Indeed, the truncation through the stated $T_{n,p}$ has
$p-1<q$ dividers, maximum property, its inherited
good hierarchy and terminal certificate (C).  By (6.1) it is elementary,
and the first consequence above converts the two full components from
the failure of interchangeability into $\mathsf{LegBad}$.  Thus the
comparison class is nonempty.  Choose a lexicographically least member,
rename its coloring $\varphi$, terminal tree $S$, and colors
$\alpha,\beta$, put $p=d+1$, and write
$A=T_{n,p-1}^{*}$.  The input coloring, corresponding ETT, hierarchy and
generating coloring are allowed to vary in this minimum; only
$G,k,e_0,n$, the lower package, (6.1), $q$, and the large estimate remain
fixed.

The chosen presentation gives two disjoint exit legs $P_1,P_2$ with
distinct outside ends.  By (6.1), $S$ is elementary.  The color $\alpha$
is missing at a unique vertex by \noderef{TreeMissingColors} and
\noderef{ColoredVertexSetElementary}, and is singleton-closed by
$\mathsf{LegBad}$.  Terminal certificate (C) says exactly
\[
 S\text{ is }
 \left(\bigcup_{\eta\in D_{n,p}}\Gamma^{p-1}_\eta\right)^-
 \text{-closed under }\varphi.                               \tag{3}
\]
in the sense of \noderef{TreeColorSetMinusClosed}.  Write $a_j$ for the
outside end and $b_j$ for the tree end of $P_j$, and order the two so
that $b_1\prec b_2$ using \noderef{TreeVertexOrder}.

We also establish once the predecessor color used in Subcases 1.2 and
2.2.  The bound \noderef{HierarchyStoredColorsCardinality} gives
$|D_{n,j}|\leq n$.  The `twoColors' and `disjointPairs' fields of
\noderef{GoodHierarchy} therefore give
\[
 \left|\bigcup_{\eta\in D_{n,j}}\Gamma^j_\eta\right|
       =2|D_{n,j}|\leq2n.                                    \tag{7}
\]
If $p=1$ and the terminal kind is not PE, (L) and (7), after deleting
$\alpha,\beta$, leave at least nine choices for
$\theta\in\overline\varphi(T_n)$ outside the terminal reserved union.
The RE/SE argument used in the non-PE base case closes every missing
color on $T_n=A$, while (3) closes $\theta$ on $S$.

If $p=1$ and the kind is PE, the `parallelBase' field of
\noderef{GoodHierarchy} gives at least four common missing colors of
$T_n$ and $R_n$ outside $\Gamma^0$.  Delete $\alpha,\beta$ and choose
$\theta$ from the at least two remaining colors.
\noderef{ParallelCommonMissingClosed}, with the pre-step elementarity
verified in the PE base case, closes $\theta$ on $A=T_n\vee R_n$.
Since $D_{n,1}\subseteq D_{n,0}$ follows directly by unfolding
\noderef{HierarchyStoredColorsAt} and using $A\subseteq S$, the choice
also avoids the reserved union in (3), which closes $\theta$ on $S$.

If $p>1$, missing-color monotonicity along the stages and (L) give at
least $2n+11$ colors missing on $T_{n,p-2}$.  Delete
$\alpha,\beta$ and the at most $2n$ colors in
$\bigcup_{\eta\in D_{n,p-1}}\Gamma^{p-2}_\eta$, and choose $\theta$.
The minus-closed field at $p-1$ closes $\theta$ on
$T_{n,p-1}=A$.  The `uniqueRefresh' field of
\noderef{GoodHierarchy} keeps every nonrefreshed pair equal to its old
pair.  If the refreshed pair contained $\theta$, its missing vertex in
$V(T_{n,p-1})-V(T_{n,p-2}^{*})$ and a $\theta$-missing vertex in
$T_{n,p-2}$ would contradict elementarity of $T_{n,p-1}$.  Thus
$\theta$ avoids the terminal reserved union and (3) closes it on $S$.
In every case $\theta\ne\alpha,\beta$ and
\[
 \begin{cases}
 \theta\in\overline\varphi(T_n),\quad
 \theta\text{ closed on }A\text{ and }S,&p=1,\\
 \theta\in\overline\varphi(T_{n,p-1}),\quad
 \theta\text{ closed on }A\text{ and }S,&p>1.
\end{cases}                                                  \tag{9}
\]

For each retained exit path, take the final edge-node before its tree
end.  It is an exact boundary identity by \noderef{TreeExitPath} and
\noderef{BoundaryEdges}.  Its color is $\alpha$ or $\beta$ by
\noderef{TwoColorIncidenceGraph}, and it cannot be $\alpha$ because the
singleton $\{\alpha\}$ is closed on $S$.  Thus these are the required
individual $\beta$-colored identities $f_1,f_2$, incident with
$b_1,b_2$ and lying on the respective path supports.

The lexicographic choice gives the two asserted comparison clauses.
Any presentation with fewer than $d$ dividers has smaller primary
coordinate, so it cannot contain two qualifying exit legs; this is the
primary no-better clause, with all of its coloring, correspondence,
maximum, large-bound, good-hierarchy, and terminal-certificate premises
exactly those in the definition of the comparison class.  If $\beta$ is
missing on $S$, the secondary coordinate is zero and the first
alternative in the theorem holds.

Otherwise, the component-by-component exit-leg construction above gives
three pairwise support-disjoint legs with distinct outside ends.  Order
their three distinct tree ends using \noderef{TreeVertexOrder}, calling
them $b_1\prec b_2\prec b_3$, and retain the corresponding ordered paths
$P_1,P_2,P_3$.  Their last boundary identities are again exact
$\beta$-colored edges $f_1,f_2,f_3$.  Choose such a triple of least total
walk length.  A same-divider comparison presentation with its second
color missing would have secondary coordinate zero; one with that color
present and a shorter qualifying triple would have a strictly smaller
positive secondary coordinate.  Either contradicts lexicographic
minimality.  This is precisely the asserted secondary no-better clause,
and the construction includes all pairwise support and endpoint
inequalities required there.

Finally take as the witnesses the chosen arbitrary input, coloring
$\varphi$, same-rung ETT, base, hierarchy and good family, followed by
$\alpha,\beta,\theta$, the ordered exit legs, and their exact final edge
identities.  The definition of the comparison class supplies properness,
one-uncoloredness, maximum property, the large estimate and (C); (6.1)
supplies elementarity; (9) supplies every field for $\theta$; and the two
preceding paragraphs supply the two minimality clauses and the final
missing-versus-three-leg alternative.  Thus every component of the
claimed global-minimum witness package has been constructed.
\end{proof}
